\subsection*{Checklist}

The syntax collected here means the following:

\begin{table}[htbp]
\centering\small
\begin{tabularx}{\linewidth}{@{}lX@{}}
\toprule
\textbf{Syntax} & \textbf{Meaning} \\
\midrule
\texttt{dma.\allowbreak{}copy.\allowbreak{}async .\allowbreak{}.\allowbreak{}.\allowbreak{}} & Start a non-blocking DMA transfer. \\
\texttt{wait f} & Wait until future \texttt{f} is ready to consume. \\
\texttt{future.data} & Access the moved tile. \\
\texttt{future.span} & Access the moved tile's shape. \\
\texttt{dma.\allowbreak{}transp\textless{}.\allowbreak{}.\allowbreak{}.\allowbreak{}\textgreater{}} & Permute dimensions during transfer. \\
\texttt{dma.\allowbreak{}pad\textless{}.\allowbreak{}.\allowbreak{}.\allowbreak{}\textgreater{}} & Pad during transfer with an explicit fill value. \\
\texttt{.zfill} & Zero-fill out-of-bounds positions in partial tiles. \\
\texttt{.\allowbreak{}swiz\textless{}N\textgreater{}} & Choose a shared-memory swizzle layout. \\
\texttt{chunkat(...)} & Select one equal tile from a regular tiling. \\
\texttt{subspan(.\allowbreak{}.\allowbreak{}.\allowbreak{}).\allowbreak{}at(.\allowbreak{}.\allowbreak{}.\allowbreak{})} & Select a tile by explicit tile extents and tile coordinates. \\
\texttt{view(.\allowbreak{}.\allowbreak{}.\allowbreak{}).\allowbreak{}from(.\allowbreak{}.\allowbreak{}.\allowbreak{})} & Select a tile by explicit extents and explicit element offsets. \\
\texttt{tile \# i} & Compose tile and in-tile coordinates into a global coordinate. \\
\texttt{\#tile} & Ask for the extent of the \texttt{tile} axis. \\
\bottomrule
\end{tabularx}
\end{table}

When reading a movement statement, first identify the selected region and its destination, then check when the destination becomes ready. These questions remain the same in a pipeline; the extra buffers simply allow different tiles to be at different points in that sequence.
